\documentclass[11pt]{article}
\usepackage[margin=1in]{geometry}
\usepackage{amsmath,amssymb,amsthm}
\usepackage[T1]{fontenc}
\usepackage{booktabs}
\usepackage{array}
\usepackage{graphicx}
\usepackage{algorithm}
\usepackage{algpseudocode}
\usepackage{hyperref}
\hypersetup{hidelinks}
\newcommand{\Oh}{\mathcal{O}}
\newcommand{\Ot}{\tilde{\mathcal{O}}}
\title{A Survey of Exact Algorithms for Subset Sum and Its Variants}
\author{Le Duc Minh\\
\small\href{mailto:minh.leduc.0210@gmail.com}{minh.leduc.0210@gmail.com}}
\date{October 2026}
\begin{document}
\maketitle
\begin{abstract}
Subset sum asks whether some of $n$ given integers add up to a target. For
every input, the best known exact algorithm still takes about $2^{n/2}$ steps,
as it did in 1974. A great deal has happened around that bound. On random
inputs, algorithms based on the representation technique run in about
$2^{0.283n}$ steps. When the numbers are small, algorithms that are nearly
linear in the target are known. For several close relatives of subset sum,
worst-case bounds below $2^{n/2}$ have been proved in the last few years,
including $2^{n/3}$ for the pigeonhole version of equal subset sum. This
survey organises these results by the question each one answers: how fast on
every input, how fast on random inputs, and how fast when the numbers are
small. For each line of work we explain the main idea, compare the bounds, and
say what the idea needs from the input. We then use these comparisons to
explain why plain subset sum has resisted so far, collect the known lower
bounds, and list open problems.
\end{abstract}
\medskip
\noindent\textbf{Keywords:} subset sum, knapsack, equal subset sum, exact
exponential algorithms, meet-in-the-middle, representation technique, survey

\section{Introduction}

\subsection{The problem}
Given integers $a_1,\dots,a_n$ and a target $b$, \emph{subset sum} asks
whether some of the $a_i$ add up exactly to $b$; in vector form, whether some
$x\in\{0,1\}^n$ satisfies $\sum_ia_ix_i=b$. The search version also asks for
such an $x$. For example, $3,5,8,13$ reach $16$ as $3+13$ and as $3+5+8$, but
never reach $2$. Subset sum is the yes/no form of the 0--1 knapsack problem and
one of the standard NP-complete problems. It also has a long history in
cryptography: the Merkle--Hellman cryptosystem~\cite{mh} based its security on
it, and the algorithms surveyed here are still used to analyse attacks on
modern schemes.

\subsection{The central question}
Since NP-complete problems are not expected to have polynomial algorithms, the
question for subset sum is how small the exponential part can be. Trying all
subsets costs $2^n$ steps. In 1974 Horowitz and Sahni~\cite{hs} brought this
down to about $2^{n/2}$, and no algorithm known today beats
$2^{(1/2-\varepsilon)n}$ on every input for any constant $\varepsilon>0$.
Whether such an algorithm exists, and in particular whether $2^{n/3}$ is
possible, is one of the best-known open problems in exact algorithms. The
constant matters in practice: for $n=100$, $2^{n/2}$ is about $10^{15}$ while
$2^{n/3}$ is about $10^{10}$.

\subsection{Three ways to measure progress}
Results on subset sum are stated in three different ways, and much of the
apparent progress on ``beating $2^{n/2}$'' depends on which one is meant
(Table~\ref{tab:regimes}).
\begin{itemize}
\item \emph{Worst case in $n$.} A bound in terms of $n$ that holds for every
input. This is the setting of the central question.
\item \emph{Average case.} A bound for inputs drawn at random, usually of
density about $1$ (numbers of about $n$ bits), often under an extra
unproved assumption about how random inputs behave.
\item \emph{Pseudopolynomial.} A bound in terms of $n$ and the target $t$ (or
the largest number), which is fast when the numbers are small.
\end{itemize}
An average-case or pseudopolynomial bound below $2^{n/2}$ says nothing about the
worst case, because the hard inputs may be exactly the ones the bound
excludes.

\begin{table}[t]
\centering
\caption{Three ways to measure progress, and the best known classical bound in
each.}
\label{tab:regimes}
\small
\begin{tabular}{p{3.2cm}p{4.2cm}p{4.6cm}}
\toprule
Setting & Best known & Where it is hard \\
\midrule
Worst case in $n$ & $2^{n/2}/n^{0.5023}$~\cite{chen} & numbers of about $n$ bits \\
Average case & $2^{0.283n}$, heuristic~\cite{bbss} & density close to $1$ \\
Pseudopolynomial & $\Ot(n+t)$~\cite{bri} & large target $t$ \\
\bottomrule
\end{tabular}
\end{table}

\subsection{Scope and organisation}
We cover exact algorithms for subset sum and its closest relatives, both
classical and quantum. We do not cover approximation schemes, the optimisation
version of knapsack, or parallel algorithms. Section~\ref{sec:prelim} fixes
notation and the variants. Sections~\ref{sec:classical}--\ref{sec:variants}
cover the worst-case, average-case and promise results in turn.
Section~\ref{sec:pseudo} covers small numbers, Section~\ref{sec:crypto} random
inputs at other densities, and Section~\ref{sec:lower} lower bounds.
Section~\ref{sec:landscape} compares all of these and draws out common themes,
and Section~\ref{sec:frontier} lists open problems.

\section{Problems and Notation}\label{sec:prelim}

\subsection{Measures}
We write $n$ for the number of integers, $t$ for the target, and
$W=\sum_i|a_i|$. The \emph{density} of an input is
$\delta=n/\log_2\max_i|a_i|$; density $1$ means the numbers have about $n$ bits.
We write $\Oh^*(f)$ for $f$ times a polynomial in $n$ and the input length, and
$\Ot(f)$ for $f$ times a polylogarithmic factor.

\subsection{One problem, many coefficient sets}
Most variants fit one template. Fix a finite set $C$ of integers, the
\emph{coefficient set}, and ask for $c\in C^n$ with $\sum_ic_ia_i$ equal to the
target, with $c\neq0$ when the target is $0$. Randolph and
W\k{e}grzycki~\cite{rw} call this \emph{subset balancing}. Subset sum is
$C=\{0,1\}$. Partition, which splits the numbers into two groups of equal sum,
is $C=\{-1,1\}$ with target $0$. Equal subset sum (ESS), which asks for two
different subsets with the same sum, is $C=\{-1,0,1\}$ with target $0$. We
write $[\pm d]=\{-d,\dots,-1,1,\dots,d\}$ and $[-d:d]=\{-d,\dots,d\}$.
Table~\ref{tab:variants} lists the variants used in this survey.

\begin{table}[t]
\centering
\caption{Problem variants, coefficient sets, and promises.}
\label{tab:variants}
\small
\begin{tabular}{p{2.9cm}p{4.4cm}}
\toprule
Variant & Coefficient set / promise \\
\midrule
Subset Sum & $C=\{0,1\}$, target $b$ \\
Partition & $C=\{-1,1\}$, target $0$ \\
Equal Subset Sum (ESS) & $C=\{-1,0,1\}$, target $0$ \\
Pigeonhole ESS (PESS) & ESS on positive $a_i$ with $\sum_ia_i<2^n-1$ (solution guaranteed) \\
Subset Balancing & fixed $C$, e.g.\ $[\pm d]$ or $[-d:d]$ \\
$k$-SUM & one element from each of $k$ lists, sum to target \\
Fixed weight & $C=\{0,1\}$ with exactly $k$ items chosen \\
Modular & the equation holds modulo some $M$ \\
\bottomrule
\end{tabular}
\end{table}

\subsection{How the variants relate}
Partition and subset sum are equally hard up to polynomial factors: Partition
is subset sum with target $W/2$, and a subset sum instance becomes a Partition
instance after adding two suitable numbers. ESS looks for a \emph{collision}
between two subset sums rather than a specific value, so its trivial bound is
$3^{n/2}$ (meet-in-the-middle over $\{-1,0,1\}$) rather than $2^{n/2}$. PESS adds
the promise $\sum_ia_i<2^n-1$. Then $2^n$ subsets share fewer than $2^n-1$ possible
sums, so by the pigeonhole principle two of them must have the same sum, and
the problem is to find them. Papadimitriou~\cite{pap} introduced PESS as a
natural problem in the complexity class PPP of such pigeonhole searches.

\section{Worst-Case Algorithms}\label{sec:classical}
This section follows the main line of worst-case results in $n$. Its story is
short: the time bound $2^{n/2}$ has stood since 1974, and most later progress
has been about memory or polynomial factors.

\subsection{Meet-in-the-middle}
Horowitz and Sahni~\cite{hs} split the numbers into a left and a right half,
list the $2^{n/2}$ subset sums of each half, sort one list, and for each left sum
$u$ search the other for $b-u$. A solution exists exactly when some pair
matches, because every subset splits uniquely into a left part and a right
part. This takes $\Oh^*(2^{n/2})$ time and memory, is deterministic, and returns
a witness. Every later worst-case algorithm in this section is a refinement of
this idea.

\subsection{Saving memory}
The $2^{n/2}$ stored sums are often the practical bottleneck. Schroeppel and
Shamir~\cite{ss} split the input into four quarters instead of two halves.
They sort the four quarter lists, of $2^{n/4}$ entries each, and use two
priority queues to produce the sums of the first two quarters in increasing
order and the sums of the last two quarters in decreasing order. A two-pointer
sweep over these streams finds a match without ever storing the half lists.
The time stays $\Oh^*(2^{n/2})$, and the memory drops to $\Oh^*(2^{n/4})$. More
generally, with memory $S\le2^{n/4}$ one can get time about $2^n/S^2$.

For about forty years $2^{n/4}$ memory was the best known at time $2^{n/2}$.
Nederlof and W\k{e}grzycki~\cite{nw} broke this barrier with a randomized
algorithm that uses memory $\Oh^*(2^{0.249999n})$, by relating subset sum to
the Orthogonal Vectors problem. Belova, Chukhin, Kulikov and
Mihajlin~\cite{bckm} lowered the memory exponent to $0.246$. These results do
not change the time.

At the other extreme, one can ask for memory polynomial in $n$. Lokshtanov and
Nederlof~\cite{ln} gave pseudopolynomial algorithms in polynomial space, and
Bansal, Garg, Nederlof and Vyas~\cite{bgnv} solved subset sum in
$\Oh^*(2^{0.86n})$ time with polynomial memory.

\subsection{Shaving polynomial factors}
The first time improvement over Horowitz and Sahni for every input came almost
fifty years later. Chen, Jin, Randolph and Servedio~\cite{chen} solved subset
sum in $\Oh(2^{n/2}/n^{\gamma})$ time for a constant $\gamma>0.5023$, in the
standard word-RAM model. Their algorithm combines representation-style ideas
with bit-packing, so that many comparisons fit into one machine word. The
exponent $1/2$ is unchanged.

\subsection{Summary}
For every input, the time exponent of subset sum has been $1/2$ for fifty years.
Memory has improved from $2^{n/2}$ to below $2^{n/4}$, and the time has improved
by a polynomial factor. The ideas that do lower the exponent come from
the average-case setting, which we turn to next.

\section{Random Inputs: The Representation Technique}\label{sec:repr}

\subsection{The key idea}
Meet-in-the-middle fixes one split of the \emph{positions} and needs a solution
to split along it. Howgrave-Graham and Joux~\cite{hgj} instead split the
\emph{solution}. A solution $x$ that picks $\ell$ numbers can be written as
$x=y+z$, where $y$ and $z$ each pick $\ell/2$ of those numbers. There are
$\binom{\ell}{\ell/2}\approx2^\ell$ ways to do this, and any one of them is enough
to find $x$. So the algorithm can afford to throw most of them away. It picks a
modulus $M$ close to the number of splits and a random residue $R$, and keeps
only the vectors $y$ whose sum is $R$ modulo $M$. If the residues of the splits
are spread evenly, about one split survives, while the list of candidates
shrinks by a factor of $M$.

The ideal two-way version gives lists of size about $2^{0.3113n}$ for a
solution of weight $n/2$. The difficulty is to build these filtered lists in
time close to their size, rather than by listing all candidates first.

\subsection{From $0.337$ to $0.283$}
Howgrave-Graham and Joux built the lists by applying the same idea one level
down: they split the solution into four pieces, match pairs of pieces under one
modulus, and then match the two pair sums exactly. Merging two lists of size
$L$ under a modulus $M$ produces about $L^2/M$ pairs, and this merge cost sets
the final time. With the top modulus $2^{\beta n}$, the time is
$\max(2^{0.3113n},2^{(0.0872+\beta)n})$ and the memory is $2^{(0.5436-\beta)n}$
for $1/4<\beta<1/3$. Table~\ref{tab:hgj} shows the resulting trade-off; the
fastest setting gives time $2^{0.337n}$.

\begin{table}[t]
\centering
\caption{Time--memory trade-off of the four-way algorithm of~\cite{hgj}, as
exponents of $2^n$.}
\label{tab:hgj}
\small
\begin{tabular}{rrl}
\toprule
$\beta$ & time & memory \\
\midrule
$\to1/4$ & 0.337 & 0.294 \\
$0.2936$ & 0.381 & 0.250 \\
$0.2972$ & 0.385 & 0.247 \\
$\to1/3$ & 0.421 & 0.211 \\
\bottomrule
\end{tabular}
\end{table}

Becker, Coron and Joux~\cite{bcj} made the pieces more flexible. They allow
the coefficient $-1$ inside the pieces, so that an unchosen number can be split
as $0=1+(-1)$ as well as $0+0$. This gives many more splits per solution, which
allows stronger filters and smaller lists. With three levels of splitting,
their time is $2^{0.291n}$. Bonnetain, Bricout, Schrottenloher and
Shen~\cite{bbss} allowed four coefficient values and reached $2^{0.283n}$, the
best classical bound for random inputs that we know of.

\subsection{The price of flexibility}
More flexible pieces also create \emph{pseudo-solutions}: pairs of pieces that
pass every filter and add up to the target, but whose sum is not a valid
$0/1$ vector, for instance because it has a $2$ or a $-1$ in some position.
The algorithm has to reject these, and if there are too many the lists grow
again. Later worst-case work handles this by checking that a pair is
\emph{compatible} on a small random sample of positions before combining
it~\cite{rw}.

\subsection{Why these bounds are average-case}
All of these analyses assume that the residues of a solution's splits are
spread evenly modulo $M$, and that the filtered lists have their expected
sizes. For random numbers of about $n$ bits this is very plausible, and
experiments agree. For a chosen input it can fail: if every number is a
multiple of $M$, every sum has residue $0$ and the filter removes nothing.
Turning the technique into a worst-case algorithm means handling such inputs,
which is the subject of the next section.

\section{From Random to Worst Case}\label{sec:rw}

\subsection{Mixing and structure}
Randolph and W\k{e}grzycki~\cite{rw} recently made the representation technique
work on every input for many coefficient sets. Their main tool is a
\emph{dichotomy}. Either the splits of a solution are spread well modulo the
filter, in which case the filter works as in the random case, or they are not,
in which case the input has special additive structure. Examples of such
structure are a small set of distinct subset sums, or exponentially many
solutions. That structure can then be used by a different algorithm, such as
subsampling or a dense dynamic program. Either way the input is solved faster
than $2^{n/2}$.

\subsection{Coefficient shifting}
The second ingredient creates extra splits when the coefficient set does not
provide them directly. They write the coefficient set as a \emph{sumset}
$C=C_1+C_2=\{c_1+c_2:c_1\in C_1,c_2\in C_2\}$ and split each coefficient vector
into one part over $C_1$ and one over $C_2$. This gives many splits only if some
coefficient has at least two representations $c_1+c_2$. For example,
$\{0,1\}+\{-3,-2,1,2\}=[\pm3]$, where $-2=0+(-2)=1+(-3)$ and $2=0+2=1+1$.

\subsection{What is solved and what is open}
With these tools they beat $|C|^{n/2}$ in the worst case for $[-d:d]$ with
$d\ge2$, for $[\pm d]$ with $d\ge3$, and for ESS, which they solve in
$1.7067^n$ time (Table~\ref{tab:rw}). Three coefficient sets remain open:
$\{0,1\}$ (subset sum), $\{-1,1\}$ (Partition) and $[\pm2]=\{-2,-1,1,2\}$. None
of the three can be written as $C_1+C_2$ so that some coefficient has two
representations. For the two-element sets this follows from
$|C_1+C_2|\ge|C_1|+|C_2|-1$. For $[\pm2]$ it follows from the fact that equality
in that inequality forces an arithmetic progression, which $[\pm2]$ is not. So
coefficient shifting has nothing to work with on exactly the cases that remain
open. This suggests that a worst-case algorithm for subset sum below $2^{n/2}$
needs a different source of extra splits.

\begin{table}[t]
\centering
\caption{Worst-case subset balancing over coefficient sets $C$~\cite{rw}. The
baseline is meet-in-the-middle, $|C|^{n/2}$.}
\label{tab:rw}
\small
\begin{tabular}{ll}
\toprule
$C$ & Worst-case result \\
\midrule
$[-d:d]$, $d\ge2$ & $|C|^{(1/2-\epsilon)n}$ for some $\epsilon>0$ \\
$[\pm d]$, $d\ge3$ & $|C|^{(1/2-\epsilon)n}$ for some $\epsilon>0$ \\
$\{-1,0,1\}$ (ESS) & $1.7067^n$ (baseline $1.7321^n$) \\
$\{-1,1\}$ (Partition) & open \\
$[\pm2]$ & open \\
$\{0,1\}$ (Subset Sum) & open \\
\bottomrule
\end{tabular}
\end{table}

\section{Where $2^{n/2}$ Has Been Beaten}\label{sec:variants}
Several close relatives of subset sum now have worst-case algorithms below
their meet-in-the-middle baseline. Each one has some feature that plain subset
sum lacks, and comparing these features shows what a breakthrough for subset
sum might need.

\subsection{Pigeonhole equal subset sum}
In PESS the numbers are positive and add up to less than $2^n-1$. There are
$2^n$ subsets but fewer than $2^n-1$ possible sums, so two subsets must share a
sum. A solution always exists, and the task is to find one.

\paragraph{A simple $2^{n/2}$ algorithm.}
The pigeonhole argument can be turned into a binary search. Call an interval
$[\ell,r]$ \emph{crowded} if more than $r-\ell+1$ subsets have their sum in it.
A crowded interval must contain two subsets with the same sum. The whole range
$[0,2^n-2]$ is crowded. Split a crowded interval in half: at least one half is
crowded, because otherwise the two halves would together hold too few subsets.
After about $n$ halvings the interval is a single crowded value, and two
subsets with that sum can be recovered. Each step needs the number of subsets
with sum at most some value $T$, which meet-in-the-middle computes in
$\Oh^*(2^{n/2})$ time, so the whole search takes $\Oh^*(2^{n/2})$.
Algorithm~\ref{alg:pess} summarises it.

\begin{algorithm}[t]
\caption{Pigeonhole binary search for PESS}
\label{alg:pess}
\begin{algorithmic}[1]
\Require positive integers $a_1,\dots,a_n$ with $\sum_ia_i<2^n-1$
\Ensure two different subsets with equal sum
\State $\ell\gets0$; $r\gets\sum_ia_i$
\While{$\ell<r$}
  \State $m\gets\lfloor(\ell+r)/2\rfloor$
  \If{$\#\{S: \ell\le w(S)\le m\}>m-\ell+1$} \Comment{count by meet-in-the-middle}
    \State $r\gets m$
  \Else
    \State $\ell\gets m+1$
  \EndIf
\EndWhile
\State \Return two subsets with sum $\ell$
\end{algorithmic}
\end{algorithm}

\paragraph{Breaking $2^{n/2}$.}
Jin and Wu~\cite{jw} gave the first faster algorithm, in $\Oh^*(2^{0.4n})$
time. Their key quantity is the number $d$ of values in $[0,2^n)$ that are
\emph{not} subset sums. This equals the total number of ``extra'' subsets that
collide with an earlier one, so it measures how many solutions there are. They
split into two cases.
\begin{itemize}
\item If $d$ is small, almost every value is a subset sum. They show that the
sorted numbers must then be very close to $1,2,4,8,\dots$: each $a_i$ differs
from $2^{i-1}$ by at most a small multiple of $d$. For such numbers, the count
of subsets with sum at most $T$ can be computed using meet-in-the-middle on only
a short prefix, so each binary-search step becomes much cheaper.
\item If $d$ is large, there are many colliding pairs. They pick a random
prime $p$ and a random residue class modulo $p$, sample subsets from that
class, and look for two with equal sum. Many collisions make this likely to
succeed quickly.
\end{itemize}
Balancing the two cases gives $2^{0.4n}$. Jin, Williams and Zhang~\cite{jwz}
then improved the bound to $\Oh^*(2^{n/3})$. Jin and Wu also gave an
$\Oh^*(2^{0.75n})$ algorithm that uses only polynomial memory.

PESS is the clearest example of a relative of subset sum that reaches
$2^{n/3}$ in the worst case. Its promise does two things that plain subset sum
lacks. It guarantees a solution, and it forces either many solutions or a very
rigid structure, and both cases can be used.

\subsection{Equal subset sum}
ESS asks for two different subsets with the same sum, with no promise. Its
meet-in-the-middle baseline is $3^{n/2}\approx1.7321^n$, since each number can
go to the first subset, the second, or neither. Mucha, Nederlof, Pawlewicz and
W\k{e}grzycki~\cite{mnpw} were the first to beat it, with $1.7088^n$, using a
worst-case form of the representation technique. The key point is that ESS has
the coefficient $0$ and the target $0$, so a solution has many representations
for free: $0=0+0=1+(-1)=(-1)+1$. Randolph and W\k{e}grzycki~\cite{rw} improved
this to $1.7067^n$, and Yamano and Shibuya~\cite{ys} to $1.6994^n$.

\subsection{Either-or subset sum}
Randolph~\cite{rand} studied a hybrid problem: given an input, return
\emph{either} a subset with the target sum \emph{or} two different subsets with
equal sum. He showed that this can be solved in randomized $\Oh^*(2^{0.461n})$
time on every input. This has a direct consequence for subset sum itself. On
an input whose $2^n$ subset sums are all different, there is no equal-sum pair
to return, so the algorithm must find a subset with the target sum whenever one
exists. Subset sum is therefore solvable below $2^{n/2}$ on all such inputs.
This is the closest known worst-case result to subset sum. It shows that any
inputs that block a faster subset sum algorithm must have some colliding
subset sums.

\subsection{Quantum algorithms}
On a quantum computer, collision-finding methods based on Grover search solve
subset sum in $\Ot(2^{n/3})$ time on every input. Bernstein, Jeffery, Lange
and Meurer~\cite{bjlm} combined quantum search with the representation
technique and obtained a heuristic average-case bound of $2^{0.241n}$. Chukhin,
Kulikov, Levitskii and Mihajlin~\cite{cklm} recently gave a worst-case quantum
bound of $\Ot(2^{2n/7})\approx2^{0.286n}$ through a worst-case quantum $k$-SUM
algorithm. So the exponent $1/3$ is not a natural limit in general; the
classical gap is a matter of algorithms we have not found yet.

\subsection{Fixed weight and $k$-SUM}
If the solution must contain exactly $k$ numbers, the problem becomes $k$-SUM:
choose $k$ numbers that add up to the target. For constant $k$,
meet-in-the-middle solves it in about $n^{\lceil k/2\rceil}$ time, and the
$k$-SUM conjecture states that this is essentially optimal. When $k$ is a
constant fraction of $n$, the representation technique applies as before.
Under the extra constraint of polynomial memory, Esser and May~\cite{em} gave a
heuristic $2^{0.65n}$ algorithm for random subset sum, improving the earlier
polynomial-memory bound of about $2^{0.72n}$.

\subsection{What these variants have in common}
Each relative that beats its baseline in the worst case has something extra
to work with. PESS has a guaranteed collision and a rigid structure when
collisions are rare. ESS has the coefficient $0$ and many free representations.
Either-or subset sum may return a collision whenever one exists. The quantum
algorithms can search faster. Plain subset sum has none of these: it may have
a single solution, there is no promise, and its coefficient set $\{0,1\}$ has
no useful factorisation.

\section{Small Numbers: Pseudopolynomial Algorithms}\label{sec:pseudo}

\subsection{From Bellman to near-linear time}
When the target $t$ is small, the reachable sums fit in a table. Bellman's
dynamic program from 1957 fills a table of the sums $0,\dots,t$ one number at a
time, in $\Oh(nt)$ steps. Packing table cells into machine words divides this
by the word size. Pisinger~\cite{pis} showed that $\Oh(n\,a_{\max})$ time is
possible, where $a_{\max}$ is the largest number, which is better when the
numbers are small compared to the target.

A newer line of work uses fast polynomial multiplication. The set of subset
sums of a list is encoded as a polynomial whose exponents are the reachable
sums. Combining two parts of the input then corresponds to multiplying their
polynomials, which takes $\Ot(t)$ time with the fast Fourier transform. For
example, the sums of $\{2\}$ and of $\{3\}$ are encoded by $1+x^2$ and $1+x^3$,
and their product $1+x^2+x^3+x^5$ encodes the sums $0,2,3,5$ of $\{2,3\}$.
Koiliaris and Xu~\cite{kx} used this to reach $\Ot(t\sqrt n)$ deterministically.
Bringmann~\cite{bri} reached $\Ot(n+t)$ with a randomized algorithm that splits
the numbers into groups by size and uses colour coding, and Jin and
Wu~\cite{jwsosa} gave a simpler $\Ot(n+t)$ algorithm based on formal power
series. Under the Strong Exponential Time Hypothesis, $t^{1-\epsilon}$ is not
possible~\cite{abhs}, so these algorithms are close to optimal in $t$.

\subsection{Faster than Bellman in every setting}
Later work refined the bound by other measures of the input.
Table~\ref{tab:pseudo} lists the main results. Bringmann and Nakos~\cite{bn}
gave an \emph{output-sensitive} algorithm, whose time depends on the number
$|S(X,t)|$ of distinct reachable sums up to $t$. Chen, Lian, Mao and
Zhang~\cite{clmz} gave $\Ot(n+\sqrt{a_{\max}t})$. Bringmann, Fischer and
Nakos~\cite{bfn} gave $\Ot(|S(X,t)|\sqrt n)$, which beats Bellman's algorithm
for every combination of parameters.

\begin{table}[t]
\centering
\caption{Pseudopolynomial algorithms for subset sum. Here $t$ is the target,
$a_{\max}$ the largest number and $|S(X,t)|$ the number of distinct subset sums
up to $t$.}
\label{tab:pseudo}
\small
\begin{tabular}{p{4.2cm}p{3.6cm}p{5.2cm}}
\toprule
Algorithm & Time & Notes \\
\midrule
Bellman (1957) & $\Oh(nt)$ & table of reachable sums \\
Pisinger~\cite{pis} & $\Oh(n\,a_{\max})$ & good for small numbers \\
Koiliaris--Xu~\cite{kx} & $\Ot(t\sqrt n)$ & deterministic, FFT \\
Bringmann~\cite{bri} & $\Ot(n+t)$ & randomized, colour coding \\
Jin--Wu~\cite{jwsosa} & $\Ot(n+t)$ & randomized, power series \\
Bringmann--Nakos~\cite{bn} & $\Ot(|S(X,t)|^{4/3})$ & output-sensitive \\
Chen--Lian--Mao--Zhang~\cite{clmz} & $\Ot(n+\sqrt{a_{\max}t})$ & \\
Bringmann--Fischer--Nakos~\cite{bfn} & $\Ot(|S(X,t)|\sqrt n)$ & beats Bellman everywhere \\
\bottomrule
\end{tabular}
\end{table}

\subsection{Dense inputs and additive combinatorics}
When there are many numbers compared to their size, a different phenomenon
takes over. Results in additive combinatorics, such as those of
Szemer\'edi and Vu~\cite{sv}, show that the subset sums of a large enough set
of integers contain long arithmetic progressions. Once such a progression is
found, every target inside a wide range can be reached by adjusting a few
choices. Bringmann and Wellnitz~\cite{bw} used this to decide dense instances
in near-linear time, and Chen, Mao and Zhang~\cite{cmz} proved stronger
structure theorems that also give a witness in near-linear time.

\subsection{What this means for the worst case}
Any of these algorithms beats $2^{n/3}$ when $t\le2^{n/3}$. But the hard
inputs for the central question have numbers of about $n$ bits, so $t$ can be
close to $2^n$, and pseudopolynomial bounds are then worse than
meet-in-the-middle. The dense results point the same way from the other side.
Inputs with many numbers relative to their size have structure that forces
many subset sums, and that structure makes them easy. Hardness lives where the
numbers are large and the subset sums are spread out.

\section{Random Inputs at Other Densities}\label{sec:crypto}
The representation technique of Section~\ref{sec:repr} targets density about
$1$. Random inputs of much lower or much higher density are easier, but for
different reasons and with different tools.

\subsection{Low density: lattices}
When the numbers have many more than $n$ bits, solutions are rare, and the
unique solution can often be found as a short vector in a lattice built from
the input. Lagarias and Odlyzko~\cite{lo} showed that almost all random
instances of density below $0.6463$ can be solved in polynomial time given an
oracle for the shortest vector problem. Coster et al.~\cite{cjloss} raised the
threshold to $0.9408$. In practice, lattice reduction algorithms such as
LLL~\cite{lll} and its stronger variants approximate this oracle well in low
dimension. These results are average-case, and they say nothing about density
$1$ or about the worst case.

\subsection{High density: generalised birthday}
When the numbers are small compared to $2^n$ but the problem is still posed
modulo a large number, there are many solutions, and the generalised birthday
algorithm of Wagner~\cite{wag} applies. It splits the input into $k$ lists and
merges them in a tree, at each level keeping only pairs whose sum is zero in
some low-order bits. For $k=4$ this takes about $2^{n/3}$ time and memory, and
for general $k$ about $2^{n/(1+\log_2k)}$. Lyubashevsky~\cite{lyu} and
Shallue~\cite{sha} applied this to random modular subset sum with a small
modulus and obtained subexponential bounds. These algorithms need many
solutions, so they do not apply to worst-case inputs with one solution.

\subsection{The density picture}
Table~\ref{tab:density} puts the three regimes side by side. The hardest
random inputs sit at density close to $1$. That is also where the worst-case
question is decided, because inputs of other densities can be handled by the
methods above.

\begin{table}[t]
\centering
\caption{Random subset sum by density $\delta=n/\log_2\max_ia_i$.}
\label{tab:density}
\small
\begin{tabular}{p{3cm}p{4cm}p{6cm}}
\toprule
Density & Typical method & Behaviour \\
\midrule
$\delta<0.94$ & lattice reduction~\cite{lo,cjloss} & unique solution found as a short vector \\
$\delta\approx1$ & representation technique~\cite{bcj,bbss} & about $2^{0.283n}$, the hardest case \\
$\delta\gg1$ & dynamic programming, generalised birthday~\cite{wag} & many solutions, small numbers \\
\bottomrule
\end{tabular}
\end{table}

\section{Lower Bounds and Barriers}\label{sec:lower}
No unconditional exponential lower bound is known for subset sum, since such a
bound would separate P from NP. The known lower bounds are conditional: they
hold if a widely believed hypothesis holds. Table~\ref{tab:lower} lists them.

\begin{table}[t]
\centering
\caption{Conditional lower bounds and their assumptions.}
\label{tab:lower}
\small
\begin{tabular}{p{6cm}p{7cm}}
\toprule
Bound & Assumption \\
\midrule
No $2^{o(n)}$ algorithm & Exponential Time Hypothesis~\cite{ip} \\
No $t^{1-\epsilon}2^{o(n)}$ algorithm & Strong Exponential Time Hypothesis~\cite{abhs} \\
No $n^{\lceil k/2\rceil-\epsilon}$ algorithm for $k$-SUM & $k$-SUM conjecture \\
No $2^{(1/2-\epsilon)n}$ algorithm for subset sum & meet-in-the-middle conjecture (unproved) \\
\bottomrule
\end{tabular}
\end{table}

The Exponential Time Hypothesis~\cite{ip} implies that subset sum has no
$2^{o(n)}$ algorithm, so some exponential dependence on $n$ is necessary. The
Strong Exponential Time Hypothesis implies that the pseudopolynomial bound
$\Ot(n+t)$ cannot be improved to $t^{1-\epsilon}2^{o(n)}$~\cite{abhs}. Neither
hypothesis says anything about the constant in the exponent between $0$ and
$1/2$. The $k$-SUM conjecture concerns a fixed number of chosen items and does
not transfer to the setting where about $n/2$ items are chosen.

The only statement that would rule out a $2^{n/3}$ algorithm is the
\emph{meet-in-the-middle conjecture}, which asserts that $2^{(1/2-\epsilon)n}$
is impossible. It is not implied by any standard hypothesis, and the recent
results of Section~\ref{sec:variants}, especially either-or subset sum, have
weakened the case for it: a natural problem very close to subset sum does beat
$2^{n/2}$ in the worst case. So, as far as is known, a $2^{n/3}$ algorithm for
subset sum is possible.

\section{Comparison and Common Themes}\label{sec:landscape}

\subsection{The results side by side}
Table~\ref{tab:landscape} collects the main results discussed above, with
their setting. Reading down the ``Setting'' column shows the pattern of this
survey: every bound below $2^{n/2}$ for subset sum itself is either
average-case, quantum, or for a relative with extra structure.

\begin{table}[t]
\centering
\caption{Main results. WC: worst case; avg: average case, often heuristic;
promise: holds under the PESS promise.}
\label{tab:landscape}
\small
\begin{tabular}{p{4.5cm}p{3.4cm}p{2.6cm}p{3.4cm}}
\toprule
Result & Time & Memory & Setting \\
\midrule
Horowitz--Sahni~\cite{hs} & $2^{n/2}$ & $2^{n/2}$ & WC \\
Schroeppel--Shamir~\cite{ss} & $2^{n/2}$ & $2^{n/4}$ & WC \\
Nederlof--W\k{e}grzycki~\cite{nw} & $2^{n/2}$ & $2^{0.249999n}$ & WC, randomized \\
Belova et al.~\cite{bckm} & $2^{n/2}$ & $2^{0.246n}$ & WC, randomized \\
Bansal et al.~\cite{bgnv} & $2^{0.86n}$ & polynomial & WC \\
Chen et al.~\cite{chen} & $2^{n/2}/n^{0.5023}$ & & WC, word RAM \\
Howgrave-Graham--Joux~\cite{hgj} & $2^{0.337n}$ & $2^{0.294n}$ & avg \\
Becker--Coron--Joux~\cite{bcj} & $2^{0.291n}$ & & avg \\
Bonnetain et al.~\cite{bbss} & $2^{0.283n}$ & & avg \\
Esser--May~\cite{em} & $2^{0.65n}$ & polynomial & avg \\
Either-or subset sum~\cite{rand} & $2^{0.461n}$ & & WC, relative \\
ESS, Yamano--Shibuya~\cite{ys} & $1.6994^n$ & & WC, relative \\
PESS, Jin--Williams--Zhang~\cite{jwz} & $2^{n/3}$ & & promise \\
Quantum, Chukhin et al.~\cite{cklm} & $2^{2n/7}$ & & WC, quantum \\
\bottomrule
\end{tabular}
\end{table}

\subsection{Theme 1: what to split}
All exponential algorithms in this survey split something. Meet-in-the-middle
and Schroeppel--Shamir split the \emph{positions} into fixed groups. This is
safe for every input, but it gives each solution only one way to be found. The
representation technique splits the \emph{solution} itself and gets many ways
to find it, which is what allows the filter to cut the search. The gain
depends on how many ways there are, which is why richer coefficient sets
($\{-1,0,1\}$ and beyond) give better exponents.

\subsection{Theme 2: many solutions help}
Many of the faster algorithms gain from having many solutions. Wagner's
algorithm and the dense algorithms need them. The large-$d$ case of PESS
samples from them. ESS and either-or subset sum benefit from the fact that
collisions among subset sums are common. The hard case for subset sum is the
opposite: a single solution, and too few collisions among the subset sums to
be useful.

\subsection{Theme 3: structure or randomness}
The recent worst-case results share one pattern. Either the input behaves like
a random one in the relevant sense, and the random-case method applies, or it
does not, and the failure reveals structure that another method can use. Jin
and Wu's PESS algorithm (few non-sums force near powers of two) and the work
of Randolph and W\k{e}grzycki (poor mixing forces additive structure) both
follow it. For subset sum, the open problem can be phrased in this language:
find a dichotomy whose two sides both beat $2^{n/2}$ when the coefficient set is
$\{0,1\}$.

\subsection{Theme 4: where hardness lives}
Taken together, the regimes in this survey narrow down the hard inputs. They
have numbers of about $n$ bits, so neither lattices nor dynamic programming
apply. Their subset sums cannot all be distinct, because either-or subset sum
handles that case below $2^{n/2}$. But collisions cannot be so plentiful that
sampling or dense structure theorems find them easily. And because their
coefficient set is $\{0,1\}$, coefficient shifting does not apply. Any
$2^{(1/2-\epsilon)n}$ algorithm for subset sum has to handle this middle
ground: inputs with some, but not many, colliding subset sums.

\subsection{Chronology}
Table~\ref{tab:chron} lists the milestones in order. Progress has come in
bursts: the classical algorithms of the 1970s and 1980s, the representation
technique around 2010, and a cluster of worst-case results for relatives since
2019.

\begin{table}[t]
\centering
\caption{Selected milestones.}
\label{tab:chron}
\small
\begin{tabular}{rp{11cm}}
\toprule
Year & Milestone \\
\midrule
1957 & Bellman's dynamic program \\
1974 & Horowitz--Sahni meet-in-the-middle, $2^{n/2}$ \\
1981 & Schroeppel--Shamir, memory $2^{n/4}$ \\
1985 & Lagarias--Odlyzko lattice attack on low-density instances \\
1994 & Papadimitriou introduces PESS and the class PPP \\
2002 & Wagner's generalised birthday algorithm \\
2010 & Howgrave-Graham--Joux representation technique, $2^{0.337n}$ average case \\
2011 & Becker--Coron--Joux, $2^{0.291n}$ average case \\
2017 & Bringmann, $\Ot(n+t)$; Bansal et al., polynomial memory \\
2019 & Mucha et al., ESS below $3^{n/2}$ \\
2020 & Bonnetain et al., $2^{0.283n}$ average case \\
2021 & Nederlof--W\k{e}grzycki, memory below $2^{n/4}$ \\
2023 & Chen et al., $2^{n/2}/n^{0.5023}$ worst case \\
2024 & Jin--Wu, PESS in $2^{0.4n}$ \\
2025 & Jin--Williams--Zhang, PESS in $2^{n/3}$; Randolph--W\k{e}grzycki, worst-case balancing \\
2026 & Yamano--Shibuya, ESS in $1.6994^n$; Chukhin et al., quantum $2^{2n/7}$ \\
\bottomrule
\end{tabular}
\end{table}

\section{Open Problems}\label{sec:frontier}

\paragraph{1. Subset sum below $2^{n/2}$.}
Is there an algorithm that solves subset sum in $\Oh^*(2^{(1/2-\epsilon)n})$
time on every input, for some $\epsilon>0$? This is the central question, and
$2^{n/3}$ would be a natural target since it matches PESS. Partition and
balancing over $[\pm2]$ are open in the same sense.

\paragraph{2. A source of representations for $\{0,1\}$.}
The worst-case representation technique needs many representations of a
solution. For $\{0,1\}$ these cannot come from coefficient shifting. Can they
come from somewhere else, for example from random restrictions of the input,
from a change of target, or from combining several related instances?

\paragraph{3. Closing the gap between random and worst case.}
On random inputs subset sum takes about $2^{0.283n}$ steps. Is there a
worst-case algorithm that matches this up to a small loss, perhaps via a
dichotomy of the kind used for PESS and either-or subset sum? Even a worst-case
bound of $2^{0.49n}$ would be a breakthrough.

\paragraph{4. The meet-in-the-middle conjecture.}
Can the conjecture be derived from a standard hypothesis, or refuted? Either
answer would settle the status of the central question. A conditional lower
bound that matches $2^{n/2}$ would explain fifty years without progress.

\paragraph{5. Memory.}
At time $2^{n/2}$, the best memory is about $2^{0.246n}$, and with polynomial
memory the best time is $2^{0.86n}$. What is the best time--memory trade-off
for every input?

\paragraph{6. Lessons from the relatives.}
PESS, ESS and either-or subset sum all beat their baselines. Is there a
reduction that transfers part of this progress to subset sum, for instance by
showing that inputs with many colliding subset sums are easy?

\section{Conclusion}
For every input, subset sum is still solved in about $2^{n/2}$ steps, as it was
fifty years ago. Around that bound the field has moved a lot. Random inputs can
be solved in about $2^{0.283n}$ steps, small numbers in nearly linear time in
the target, and several close relatives now have worst-case algorithms below
their baselines, down to $2^{n/3}$ for PESS. Comparing these results shows what
they rely on: many solutions, a promise, extra coefficients, or a random-like
input. Plain subset sum, with a possibly unique solution, no promise and the
coefficient set $\{0,1\}$, offers none of these. A worst-case algorithm below
$2^{n/2}$ would need a new idea, and no known lower bound rules one out.

\begin{thebibliography}{99}
\bibitem{mh} R. Merkle and M. Hellman, ``Hiding information and signatures in
trapdoor knapsacks,'' \emph{IEEE Trans. Inf. Theory} 24(5):525--530, 1978.
\bibitem{pap} C. H. Papadimitriou, ``On the complexity of the parity argument
and other inefficient proofs of existence,'' \emph{J. Comput. Syst. Sci.}
48(3):498--532, 1994.
\bibitem{pis} D. Pisinger, ``Linear time algorithms for knapsack problems with
bounded weights,'' \emph{J. Algorithms} 33(1):1--14, 1999.
\bibitem{sv} E. Szemer\'edi and V. Vu, ``Finite and infinite arithmetic
progressions in sumsets,'' \emph{Ann. of Math.} 163(1):1--35, 2006.
\bibitem{lll} A. K. Lenstra, H. W. Lenstra, L. Lov\'asz, ``Factoring polynomials
with rational coefficients,'' \emph{Math. Ann.} 261:515--534, 1982.
\bibitem{ip} R. Impagliazzo and R. Paturi, ``On the complexity of $k$-SAT,''
\emph{J. Comput. Syst. Sci.} 62(2):367--375, 2001.
\bibitem{hs} E. Horowitz and S. Sahni, ``Computing partitions with applications
to the knapsack problem,'' \emph{J. ACM} 21(2):277--292, 1974.
\bibitem{ss} R. Schroeppel and A. Shamir, ``A $T=O(2^{n/2})$, $S=O(2^{n/4})$
algorithm for certain NP-complete problems,'' \emph{SIAM J. Comput.}
10(3):456--464, 1981.
\bibitem{chen} X. Chen, Y. Jin, T. Randolph, R. A. Servedio, ``Subset sum in
time $2^{n/2}/\mathrm{poly}(n)$,'' RANDOM 2023; arXiv:2301.07134.
\bibitem{nw} J. Nederlof and K. W\k{e}grzycki, ``Improving Schroeppel and
Shamir's algorithm for subset sum via orthogonal vectors,'' STOC 2021;
arXiv:2010.08576.
\bibitem{bckm} E. Belova, I. Chukhin, A. Kulikov, I. Mihajlin, ``Improved space
bounds for subset sum,'' ESA 2024; arXiv:2402.13170.
\bibitem{bgnv} N. Bansal, S. Garg, J. Nederlof, N. Vyas, ``Faster space-efficient
algorithms for subset sum, $k$-sum, and related problems,'' STOC 2017;
\emph{SIAM J. Comput.} 47(5):1755--1777, 2018.
\bibitem{ln} D. Lokshtanov and J. Nederlof, ``Saving space by algebraization,''
STOC 2010.
\bibitem{hgj} N. Howgrave-Graham and A. Joux, ``New generic algorithms for hard
knapsacks,'' EUROCRYPT 2010; IACR ePrint 2010/189.
\bibitem{bcj} A. Becker, J.-S. Coron, A. Joux, ``Improved generic algorithms for
hard knapsacks,'' EUROCRYPT 2011; IACR ePrint 2011/474.
\bibitem{bbss} X. Bonnetain, R. Bricout, A. Schrottenloher, Y. Shen, ``Improved
classical and quantum algorithms for subset-sum,'' ASIACRYPT 2020; ePrint
2020/168.
\bibitem{rw} T. Randolph and K. W\k{e}grzycki, ``Beating meet-in-the-middle for
subset balancing problems,'' arXiv:2511.10823, 2025.
\bibitem{ys} Yamano and Shibuya, ``Faster exact algorithms for Equal-Subset-Sum,''
arXiv:2607.09289, 2026.
\bibitem{jw} C. Jin and H. Wu, ``A faster algorithm for Pigeonhole Equal Sums,''
arXiv:2403.19117, 2024.
\bibitem{jwsosa} C. Jin and H. Wu, ``A simple near-linear pseudopolynomial time
randomized algorithm for subset sum,'' SOSA 2019; arXiv:1807.11597.
\bibitem{mnpw} M. Mucha, J. Nederlof, J. Pawlewicz, K. W\k{e}grzycki,
``Equal-subset-sum faster than the meet-in-the-middle,'' ESA 2019.
\bibitem{em} A. Esser and A. May, ``Low weight discrete logarithm and subset
sum in $2^{0.65n}$ with polynomial memory,'' EUROCRYPT 2020.
\bibitem{bn} K. Bringmann and V. Nakos, ``Top-$k$-convolution and the quest for
near-linear output-sensitive subset sum,'' STOC 2020.
\bibitem{lyu} V. Lyubashevsky, ``The parity problem in the presence of noise,
decoding random linear codes, and the subset sum problem,'' RANDOM 2005.
\bibitem{sha} A. Shallue, ``An improved multi-set algorithm for the dense subset
sum problem,'' ANTS 2008.
\bibitem{cjloss} M. J. Coster, A. Joux, B. A. LaMacchia, A. M. Odlyzko,
C.-P. Schnorr, J. Stern, ``Improved low-density subset sum algorithms,''
\emph{Comput. Complexity} 2(2):111--128, 1992.
\bibitem{abhs} A. Abboud, K. Bringmann, D. Hermelin, D. Shabtay, ``SETH-based
lower bounds for subset sum and bicriteria path,'' SODA 2019.
\bibitem{jwz} C. Jin, R. Williams, S. Zhang, ``New algorithms for Pigeonhole
Equal Subset Sum,'' ESA 2025, LIPIcs 351:86.
\bibitem{rand} T. Randolph, ``Exact and parameterized algorithms for subset sum
problems,'' PhD thesis, Columbia University, 2024.
\bibitem{bjlm} D. Bernstein, S. Jeffery, T. Lange, A. Meurer, ``Quantum
algorithms for the subset-sum problem,'' PQCrypto 2013.
\bibitem{cklm} I. Chukhin, A. Kulikov, V. Levitskii, I. Mihajlin, ``Improved
quantum algorithms for subset sum and $k$-SUM,'' arXiv:2608.07309, 2026.
\bibitem{wag} D. Wagner, ``A generalized birthday problem,'' CRYPTO 2002.
\bibitem{lo} J. Lagarias and A. Odlyzko, ``Solving low-density subset sum
problems,'' J. ACM 32(1):229--246, 1985.
\bibitem{kx} K. Koiliaris and C. Xu, ``A faster pseudopolynomial time algorithm
for subset sum,'' SODA 2017.
\bibitem{bri} K. Bringmann, ``A near-linear pseudopolynomial time algorithm for
subset sum,'' SODA 2017.
\bibitem{clmz} L. Chen, J. Lian, Y. Mao, G. Zhang, ``An improved pseudopolynomial
time algorithm for subset sum,'' FOCS 2024.
\bibitem{bfn} K. Bringmann, N. Fischer, V. Nakos, ``Beating Bellman's algorithm
for subset sum,'' SODA 2025; arXiv:2410.21942.
\bibitem{bw} K. Bringmann and P. Wellnitz, ``On near-linear-time algorithms for
dense subset sum,'' SODA 2021; arXiv:2010.09096.
\bibitem{cmz} L. Chen, Y. Mao, G. Zhang, ``Long arithmetic progressions in sumsets
and subset sums,'' STOC 2025; arXiv:2503.19299.
\end{thebibliography}
\end{document}
